% Irodov Part III §3.6 first original problem 288,
% user PDF page 144 (printed 135), figure 3.78. 
\electroinductionproblem{288}{135}
图3.78：一根导线弯成抛物线
$y=ax^2$，位于垂直图面的匀强磁场
$\bm B$ 中。
从 $t=0$ 开始，一根水平导电连杆
从抛物线顶点沿 $+y$ 方向以恒定
加速度 $w$ 平动，其两端始终接触抛物线，
形成闭合回路。
求感应电动势的模随连杆高度 $y$ 的关系。
\exsolution
由抛物线方程，
高度为 $y$ 时两交点横坐标为
$x_\pm=\pm\sqrt{y/a}$，
所以连杆有效长度
\[
\ell(y)=x_+-x_-=2\sqrt{\frac ya}.
\]
连杆从静止开始以恒加速度 $w$ 运动，
由运动学
$v^2=2wy$，
故
\[
v(y)=\sqrt{2wy}.
\]
因为磁场垂直回路平面，
动生电动势的模为
$|\mathcal E|=B\ell(y)v(y)$，
代入可得
\[
|\mathcal E|=
B\,2\sqrt{\frac ya}\sqrt{2wy}
=2By\sqrt{\frac{2w}{a}}.
\]
因此
\[
\boxed{|\mathcal E(y)|=
By\sqrt{\frac{8w}{a}}}.
\]
也可从回路面积推导：
\[
S(y)=\int_0^y
2\sqrt{\frac{y'}a}\,\mathrm dy'
=\frac4{3\sqrt a}y^{3/2},
\quad
|\mathcal E|=B\frac{\mathrm dS}{\mathrm dt}
=B\frac{\mathrm dS}{\mathrm dy}v.
\]
两种方法完全一致。
\excheck $a$ 在 $y=ax^2$ 中的单位为
$\mathrm{m^{-1}}$，
因此
$By\sqrt{w/a}$ 的单位为伏特。
符号由原图指定的回路正向及
$\bm B$ 指向纸内或纸外决定；
此处报告模。
